\documentclass[10pt,a4paper]{amsart}
\usepackage[margin=19mm]{geometry}
\usepackage{amsmath,amssymb,mathtools}
\usepackage[scr=rsfs,cal=euler]{mathalfa}
\usepackage{xfrac}
\usepackage[unicode,hidelinks,bookmarks=false]{hyperref}
\newcommand{\ie}{i.e.\ }
\newcommand{\cf}{cf.\ }
\newcommand{\resp}{respectively\ }
\newcommand{\Subsection}{\subsection}
\providecommand{\nobreakdash}{\nobreak\hbox{-}}
\setlength{\emergencystretch}{2em}
\multlinegap0pt
\usepackage{bm}
\usepackage{bbm}
\usepackage{dsfont}
\usepackage{xspace}
\usepackage{cancel}
\usepackage{mathtools}
\usepackage{amsthm}
\makeatletter
\@ifundefined{theorem}{\newtheorem{theorem}{Theorem}}{}
\makeatother
\makeatletter
\newcommand{\mc}[1]{\text{$\mathcal{#1}$}}
\expandafter\ifx\csname remark\endcsname\relax
\newenvironment{remark}%
{\vskip6pt%
\noindent%
{\it Remark.}}%
{\vskip6pt}
\fi
\expandafter\ifx\csname remarks\endcsname\relax
\newenvironment{remarks}%
{\vskip6pt%
\noindent%
{\it Remarks}. %
\begin{renumerate}}%
{\end{renumerate}\vskip6pt}
\fi
\expandafter\ifx\csname renumerate\endcsname\relax
\newenvironment{renumerate}%
{\renewcommand{\labelenumi}{{\it \roman{enumi}}\hskip2pt)}%
\begin{enumerate}}%
\fi
\renewcommand{\to}{\longrightarrow}
\makeatother
\title[T-duality for transgressive fibrations]{T-duality for transgressive fibrations}
\author{Gil R. Cavalcanti}
\date{Source: arXiv:2503.20143v2 (CC BY 4.0)}
\begin{document}
\maketitle

\noindent\emph{Source and licence.} T-duality for transgressive fibrations. Version arXiv:2503.20143v2 (CC BY 4.0), distributed under the Creative Commons Attribution 4.0 International licence, as recorded in the arXiv licence metadata for this exact version and shown on the abstract page https://arxiv.org/abs/2503.20143v2. Full rights notice: licenses/arxiv-2503.20143-CC-BY-4.0.txt.\par\smallskip

\noindent\textit{Editorial scope.} Counted unit: the Theorem whose source label is \texttt{eq:multidegree H condition} in the article (source0.tex lines 526--532, source label \texttt{eq:multidegree H condition}), delivered together with the article's own complete original proof of it (source0.tex lines 533--570, one \texttt{proof} environment of 38 lines). No mathematics is written, repaired or re-derived: statement and proof are the article's own text line for line. Required context retained uncounted: none. Every definition, display and label of those lines is retained unchanged. Omitted from this package and not redistributed: 1--525, 571--712. Nothing inside a retained excerpt is omitted or altered: every retained body is delivered exactly as the article's lines are written, so no omission receipt of any kind accompanies this package. Citations: the retained excerpts carry no citation command at all, so no bibliography entry of the article is transcribed and no reference is redistributed. Reference adaptation: none was needed and none was made; every reference inside the delivered unit resolves inside it, so no unresolved reference or citation is produced. Printed slips kept literal: no printed word, symbol, hypothesis or number of the counted statement or of its proof was altered. Layout and class: the article's own class is \texttt{article} with packages including latexsym, amssymb, amsthm, amsmath, xypic, lscape, epsfig, setspace, tikz, slashed, IEEEtrantools, hyperref, caption, color; the wrapper is the lane's pinned \texttt{amsart} editorial wrapper at 10pt on A4, which restates the theorem-like environments the retained text needs and adds the article's own macro definitions that the retained text needs (\texttt{\textbackslash mc}, \texttt{\textbackslash to}), with extra wrapper packages (bm, bbm, dsfont, xspace, cancel, mathtools). The delivered numbering is the unit's own. Assets: no retained span contains a figure environment, a graphics-inclusion command, a PDF-page inclusion, a table or a TikZ picture; no image file is copied into this package, the package declares no asset dependency and no image file is redistributed with it. Licence: version arXiv:2503.20143v2 of this article is distributed under the Creative Commons Attribution 4.0 International licence, as recorded in the arXiv licence metadata for this exact version and shown on the abstract page https://arxiv.org/abs/2503.20143v2. A full rights notice is delivered at licenses/arxiv-2503.20143-CC-BY-4.0.txt. Characters: every retained excerpt is pure ASCII, so the delivered engine, XeLaTeX with its default Latin Modern text font, reports no missing character.
\begin{theorem}
A T-dual of $S^{2n-1}{\cdots} (E,H) \stackrel{\pi}{\longrightarrow} M$ exists if and only if $\pi_*H$ is the product of an integral class of fixed degree by an invertible element in $H^\bullet(M)$, that is,
\begin{equation}\label{eq:multidegree H condition}
[\pi_*H] = \hat{\mc{E}} \mc{B},
\end{equation}
where $\hat{\mc{E}}$ is an integral class and $\mc{B}$ is a mutidegree even cohomology class with nonvanishing zero degree term, $\mc{B}_0 \neq 0$.
\end{theorem}

\begin{proof}
For the necessity, if $\hat{E}$ is T-dual to $E$,  then
\begin{equation}\label{eq:restriction on H}
0 = [\hat{p}_*(dF)] = \hat{p}_*[p^*H - \hat{p}^*\hat{H}] =  \hat{p}_*\circ p^*[H] =   \hat{\pi}^*\circ \pi_*[H].
\end{equation}
Since by the Gysin sequence for $\hat{E}$, the cohomology classes of $M$ which pull back to trivial classes on $\hat{E}$ are multiples of the Euler class of $\hat{E}$ we conclude that if $\hat{E}$ exists,  \eqref{eq:multidegree H condition} holds.


If the Euler class of $\hat{E}$ vanishes, that is $\hat{\mc{E}} =0$, then we can take $\mc{B}=1$ and the claim about the decomposition of $\pi_*H$ holds.

To complete the proof that \eqref{eq:multidegree H condition} is a necessary condition, we need to consider the case when $\hat{\mc{E}} \neq 0$. Let $\psi$ be a global angular form for $E$ so that $d\psi = \pi^*e$ is a representative of the Euler class of $E$ and similarly for $\hat\psi$ on $\hat{E}$. In this case, from the Gysin sequence for the cohomology of $E\times_M \hat{E}$, there is no closed form in the correspondence space which integrates nonzero along the fibers of $E\times_M \hat{E} \to M$.
Therefore, if $F'$ is another form for which $dF' = p^*H -\hat{p}^*\hat{H}$, then $F-F'$ is closed and hence $(\pi_* \circ p_* (F-F'))_0 = 0$. Hence $(\pi_* \circ p_* (F))_0 \neq 0$ for any $F$ such that $dF =  p^*H -\hat{p}^*\hat{H}$. 

Changing $H$ and $H'$ by exact forms we can arrange that they lie in the corresponding Gysin complexes for $E$ and $\hat{E}$ and therefore we can also pick $F$ in the Gysin complex for $E\times_M \hat{E}$. So we have
\begin{equation}\label{eq:H looks like}
H = H_{[1]} +\psi  H_{[0]},
\end{equation}
\begin{equation}\label{eq:F looks like}
F = F_{[0]} +  \psi F_{[1]} +\hat\psi \hat{F}_{[1]} + \psi \hat\psi F_{[2]} ,
\end{equation}
where $H_{[i]}$ and $F_{[i]}$ are multidegree forms pulled back from $M$ whose parity agrees with the partiy of their indices and $(F_{[2]})_0 \neq 0$. Then
\[\hat{\pi}^*H_{[0]} = \hat{p}_* p^* H =\hat{p}_* (p^*H -\hat{p}^* \hat{H}) = \hat{p}_* dF = dF_{[1]} - \hat{e} F_{[2]} + \hat{\psi} dF_{[2]}. \]
Therefore $dF_{[2]}=0$ (in particular $(F_{[2]})_0$ is constant and nonzero) and since both sides are pulled back form $M$, we conclude that on $M$
\[H_{[0]}  = dF_{[1]} - \hat{e} F_{[2]}.\]
Passing to cohomology, we have that $\mc{B} =- [F_{[2]}]$ satisfies the conditions of the theorem. 

For the converse we assume \eqref{eq:multidegree H condition} holds. Let $\psi$ be a global angular form for $E$ so that $d\psi = \pi^*e$ is a representative of the Euler class of $E$. Let $\hat{E}$ be a sphere bundle whose Euler class is a constant multiple of $\hat{\mc{E}}$ (and after scaling and renaming we may assume it to be $\hat{\mc{E}}$).

Let $\hat{\psi}$ be a global angular form for $\hat{E}$ with $d\hat{\psi} = \hat{e}$, a representative for $\hat{\mc{E}}$. Since the inclusion of the Gysin complex is a quasi-isomorphism, after changing $H$ by an exact term (which does not afect the existence of T-duals) we can assume that $H \in \Omega^\bullet_\psi(E)$ and write
\[H = \pi^*H_ {[0]} + \psi \pi^*H_{[1]}.\]
With $H_{[0]}$ and $H_{[1]}$ multidegree forms. Since $[H_{[1]}]$ is a multiple of $[\hat{e}]$, after changing by another exact element we can arrange that
\[H = H_{[0]} + \psi \hat{e} H_{[1]}',\]
Let $\hat{H} = H_{[0]} + \hat{\psi} e H_{[1]}'$, then
\[H - \hat{H} = d(-\psi\hat{\psi} H_{[1]}')\]
We take
\[F = -\psi\hat{\psi} H_{[1]}'\]
 and observe that $[\pi_*\circ p_* F] = [H_{[1]}'] = \mc{B}$. In particular $[(\pi_*\circ p_* F)_0]=\mc{B}_0 \neq 0$.
\end{proof}

\end{document}
